\worksheettitle
  {Дерево выбора}
  {\modulemark{M05-R} Коррекционная карточка}

\studentnameblock

\begin{notebox}[Задача]
  Из цифр \(0,\ 3,\ 6\) составьте все двузначные числа без повторений.
  Запись двузначного числа не может начинаться с нуля.
\end{notebox}

\begin{center}
\begin{tikzpicture}[
  level distance=19mm,
  level 1/.style={sibling distance=48mm},
  level 2/.style={sibling distance=22mm},
  every node/.style={draw=Primary,rounded corners,inner sep=4pt,font=\sffamily},
  edge from parent/.style={draw=Secondary,-{Stealth}}
]
\node {первая цифра}
  child {node[draw=Accent,text=Accent] {0: нельзя}}
  child {node {3}
    child {node {0}}
    child {node {6}}}
  child {node {6}
    child {node {0}}
    child {node {3}}};
\end{tikzpicture}
\end{center}

\begin{practicebox}[Пройдите по ветвям]
  Начинаются с 3: \answerblank[18mm], \answerblank[18mm].

  Начинаются с 6: \answerblank[18mm], \answerblank[18mm].

  Полный список: \answerblank[78mm]. Всего чисел: \answerblank[13mm].
\end{practicebox}

\begin{reflectionbox}[Почему список полный]
  Ноль исключён только из \answerblank[24mm] позиции. Проверены все допустимые
  первые цифры. Для каждой проверены все разрешённые вторые цифры.
\end{reflectionbox}

\newpage
\section{Различаем условия}

\begin{practicebox}[1. Цифры 3 и 7]
  \begin{enumerate}[label=\textbf{\arabic*.}]
    \item Двузначные числа; повторения разрешены; использовать обе цифры
      необязательно:
      \answerblank[68mm].
    \item Двузначные числа без повторений: \answerblank[48mm].
    \item Трёхзначные числа; повторения разрешены; обе цифры обязательны:
      \answerblank[84mm].
  \end{enumerate}
\end{practicebox}

\begin{warningbox}[2. Найдите повтор и пропуск]
  Для цифр \(1,\ 3,\ 6\) без повторений записали:
  \(13,\ 16,\ 31,\ 36,\ 61,\ 61\).

  Повтор: \answerblank[15mm]. Пропущено: \answerblank[15mm].
  Как вы проверили список?
  \writinglines{2}
\end{warningbox}

\begin{checkpointbox}[Повторная проверка]
  \begin{enumerate}[label=\textbf{\arabic*.}]
    \item Из цифр \(0,\ 4,\ 8\) составьте все двузначные числа без повторений:
      \answerblank[70mm].
    \item Из цифр \(2,\ 5\) составьте все трёхзначные числа с повторениями,
      если обе цифры обязательны:
      \answerblank[88mm].
  \end{enumerate}
\end{checkpointbox}
